% SPDX-FileCopyrightText: © 2026 Idiap Research Institute <contact@idiap.ch>
%
% SPDX-FileContributor: Andre Freitas <andre.freitas@idiap.ch>
% Neuro-symbolic AI Group
%
% SPDX-FileContributor: Danilo Gusicuma <danilo.gusicuma@idiap.ch>
% Neuro-symbolic AI Group
%
% SPDX-License-Identifier: MIT
% Generated from step-explanations.json
\paragraph{Step-by-step reconstruction.}
A finite-result draft requests a sign certificate for the matrix with diagonal entries 1 and off-diagonal entries t, for absolute t at most one half. A geometric application still needs a separate transfer estimate.
The following is a source-guided illustrative reconstruction. Internal procedures are prescriptions; the native declarations do not establish their execution.
\begin{enumerate}[leftmargin=*,label=\arabic*.]
\item \textbf{Recover the actual expansion need} (\texttt{e0}).
\begin{description}[style=standard,labelindent=0pt,leftmargin=5.5em,labelwidth=5em,labelsep=0.5em,align=left,font=\normalfont\bfseries,itemsep=0.15em,parsep=0pt,topsep=0.3em]
\item[Purpose] Distinguish mathematical gaps from prose gaps.
\item[Input] Draft; stated contribution; expansion markers
\item[Method] Classify each request as a definition, coefficient, finite check, transfer theorem or explanatory paragraph and recover the smallest honest scope.
\item[Check] Do not treat all expansion markers as requests to assert stronger results.
\item[Output] Expansion scope: finite matrix certificate; separate application obligation
\item[Next] Extension formation consumes this scope; unresolved geometric work remains distinct.
\item[Example] Specify the two-by-two symmetric matrix, the range absolute t at most one half, and the still-unproved geometric transfer estimate.
\end{description}
\item \textbf{Propose a bounded extension} (\texttt{a0}).
\begin{description}[style=standard,labelindent=0pt,leftmargin=5.5em,labelwidth=5em,labelsep=0.5em,align=left,font=\normalfont\bfseries,itemsep=0.15em,parsep=0pt,topsep=0.3em]
\item[Purpose] Construct the missing finite certificate.
\item[Input] Matrix scope; parameter range; missing sign certificate
\item[Method] Define the carrier and entries, order coefficient extraction before the sign test, and state the lemma that will consume it.
\item[Check] A finite extension must expose its quantitative range and avoid assuming the application.
\item[Output] Extension hypothesis: exact eigenvalue-based positivity certificate
\item[Next] Refutation tests this proposed extension and its parameter boundary.
\item[Example] For the matrix with diagonal 1 and off-diagonal t, propose eigenvalues 1+t and 1-t and hence a lower bound of one half on the stated range.
\end{description}
\item \textbf{Test the extension's validity} (\texttt{r0}).
\begin{description}[style=standard,labelindent=0pt,leftmargin=5.5em,labelwidth=5em,labelsep=0.5em,align=left,font=\normalfont\bfseries,itemsep=0.15em,parsep=0pt,topsep=0.3em]
\item[Purpose] Test the finite result and its application boundary.
\item[Input] Proposed matrix certificate; range; intended use
\item[Method] Check the algebra and boundary values, seek a non-example, and localize any absent transfer lemma.
\item[Check] Finite sign success does not discharge a geometric comparison hypothesis.
\item[Output] Validity assessment: scoped matrix calculation; open transfer debt
\item[Next] Readiness joins this assessment with the same extension hypothesis.
\item[Example] At t=1/2 the eigenvalues are 3/2 and 1/2; at t=2 they are 3 and -1. The range matters, and neither calculation supplies geometric transfer.
\end{description}
\item \textbf{Integrate with its evidence boundary} (\texttt{e1}).
\begin{description}[style=standard,labelindent=0pt,leftmargin=5.5em,labelwidth=5em,labelsep=0.5em,align=left,font=\normalfont\bfseries,itemsep=0.15em,parsep=0pt,topsep=0.3em]
\item[Purpose] Assess inclusion while preserving proof obligations.
\item[Input] Extension hypothesis + its validity assessment
\item[Method] Explain how the finite lemma fits the argument, retain the counterexample and separate paper readiness from unproved application claims.
\item[Check] An integrated paragraph cannot turn a missing bridge into an accepted theorem.
\item[Output] Finite certificate; expansion validation; readiness decision
\item[Next] Original proof-required routing needs a ProofQueueItem adapter; native return does not execute it.
\item[Example] Add the exact finite positivity result and its range to the draft, while sending the geometric transfer estimate to further proof work instead of claiming the application.
\end{description}
\end{enumerate}
\paragraph{Changed-input contrast.}
If t is allowed to exceed 1, the smaller eigenvalue becomes negative and the proposed positivity extension fails. If the finite range is retained, positivity still cannot establish the missing geometric transfer.
